% Imporditud: tools/import_eksperimendid.py
% Allikas: Eesti füüsikaolümpiaadi eksperimendi ülesannete
%   kogu (Rael Kalda, 2023), ülesanne Ü101, lahendus L101.
\setAuthor{EFO žürii}
\setRound{piirkonnavoor}
\setYear{2011}
\setNumber{G E2}
\setDifficulty{3}
\setTopic{dünaamika}
\setVahendid{niit, tundmatu massiga koormis, etteantud massiga joonlaud}
\setPunktid{12}
\setAllikas{Eksperimendi kogu Ü101, lahendus lk 78}
\setLahendusseis{toores}

\prob{Niit}
Leida niidi katkemispinge (st minimaalne jõud, mille rakendamisel niidi otstele niit katkeb).

\hint

\solu
Lahendusidee on kasutada kangimeetodit niidi katkemispinge määramiseks: asetame joonlaua ühele otsale klotsi ja teise otsa külge seome niidi; kinnitamiseks kasutatav aas niidi küljes peab olema piisavalt suur, et niit ei katkeks mitte aasa kohalt vaid sõlme kohalt või (eelistatult) allpool sõlme. Kangkaalu toetuspunktina kasutame laua serva. Muutes üle lauaserva ulatuva joonlauaosa pikkust leiame maksimaalse pikkuse $l_1$, mille puhul niit katkeb, kui allapoole sikutades üritame klotsi kergitada. Selle meetodi eest 2 punkti. Kui kaugus lauaservast klotsi keskpunktini on $l_2$ ning kaugus joonlaua keskpunktini $L_1$, klotsi mass M ja joonlaua mass m, siis jõumomentide tasakaalutingimusest saame avaldada niidi katkemispinge T = (M $l_2$ + mL1)g/$l_1$ (selle valemi eest 1,5 punkti; kui puudub joonlaua massi arvestav liige, siis 1 punkt). $l_1$, $l_2$ ja $L_1$ mõõtmise eest 3 $\times$ 0,5 punkti (kui on mõõdetud teisi pikkusi, mille põhjal saab mainitud kolm pikkust arvutada, siis antakse loomulikult ikkagi 1,5 punkti). Klotsi massi leidmiseks kasutame analoogset kangkaalu meetodit: tasakaalustame joonlaua otsa paigutatud klotsi joonlaua massi abil (selle idee eest 2 punkti). Kui joonlaua keskpunkt asub lauaservast kaugusel $L_2$ ja klotsi keskpunkt kaugusel $l_3$, siis M = mL$^2$/$l_3$ (valemi eest 1 punkt, arvutuste eest 1 punkt). $L_2$ ja $l_3$ mõõtmiste eest 2 $\times$ 0,5 punkti. Korduvmõõtmiste sooritamise eest 1 punkt. Mõistlikus suurusjärgus vastuse eest 1 punkt; kui vastus erineb õigest väärtusest enam kui 1,5 korda või lõppvastusest on ühikud puudu (kuid lahenduse põhjal on arusaadav, mis ühikud peaksid olema, siis 0,5 punkti); kui vastus erineb õigest väärtusest enam kui 2 korda (või ühikud puuduvad ning lahenduse/arvutuste põhjal pole võimalik aru saada, mis ühikud peaksid olema), siis 0 punkti.
\probend
